% Problem dossier: VI.05.P5
\subsection*{Problem VI.05.P5 --- Connectedness, idempotents and product rings}
\label{prob:vi-05-p5}

\paragraph{Problem.}
\label{prob:vi05-p5}
Assume \(k\) is algebraically closed and \(X\) is an affine algebraic set.  Prove that the
following are equivalent:
\begin{enumerate}
  \item \(X\) is disconnected;
  \item \(k[X]\) contains a nontrivial idempotent;
  \item
        \[
        k[X]\cong A_1\times A_2
        \]
        for nonzero \(k\)-algebras \(A_1,A_2\).
\end{enumerate}

Then apply the criterion to:
\[
V(xy)\subseteq\mathbb A_k^2
\]
and
\[
V(x(x-1))\subseteq\mathbb A_k^2.
\]

\ifdefined\FullProblemDossiers

\paragraph{Why this problem matters.}
This problem separates two notions that are often confused: reducibility and
disconnectedness.  It gives the algebraic signature of a topological separation - a
nontrivial idempotent and a product decomposition.

\paragraph{Useful ideas before starting.}
\begin{itemize}
\item \(e^2=e\) implies \(e(e-1)=0\);
\item the Chinese remainder theorem;
\item the Nullstellensatz for translating disjoint algebraic sets into comaximal ideals.
\end{itemize}

\paragraph{Guided hints.}
\textbf{Hint 1.} A nontrivial idempotent \(e\) gives two disjoint clopen sets
\(V_X(e)\) and \(V_X(e-1)\).

\textbf{Hint 2.} A product \(A_1\times A_2\) contains the idempotent \((1,0)\).

\textbf{Hint 3.} If \(X=Y\sqcup Z\) is a separation into algebraic subsets, then
\(I(Y)+I(Z)=(1)\) over an algebraically closed field.  Use the Chinese remainder theorem.

\paragraph{Solution.}
Suppose first that \(e\in k[X]\) is a nontrivial idempotent.  Since
\[
e^2=e,
\]
at every point \(x\in X\),
\[
e(x)\in\{0,1\}.
\]
Thus
\[
X=V_X(e)\sqcup V_X(e-1).
\]
Both sets are closed and complementary, hence open.  Nontriviality of \(e\) makes both
nonempty, so \(X\) is disconnected.

Conversely, suppose
\[
X=Y\sqcup Z
\]
with \(Y,Z\) nonempty clopen.  Since both are closed algebraic subsets,
\[
Y\cap Z=\varnothing
\]
implies
\[
V(I(Y)+I(Z))=\varnothing.
\]
Over algebraically closed \(k\), the Nullstellensatz gives
\[
I(Y)+I(Z)=k[x_1,\dots,x_n].
\]
Choose
\[
f\in I(Y),
\qquad
g\in I(Z)
\]
with
\[
f+g=1.
\]
The class of \(g\) in \(k[X]\) is \(1\) on \(Y\) and \(0\) on \(Z\), hence is a
nontrivial idempotent.

If \(A\) has a nontrivial idempotent \(e\), then
\[
A\to eA\times(1-e)A,
\qquad
a\mapsto(ea,(1-e)a)
\]
is an isomorphism.  Conversely, \((1,0)\) is a nontrivial idempotent in every nontrivial
product \(A_1\times A_2\).

For
\[
V(xy),
\]
the two irreducible components meet, so the space is connected and the coordinate ring has
no nontrivial idempotent.

For
\[
V(x(x-1)),
\]
the two components are disjoint and
\[
k[x,y]/(x(x-1))
\cong
k[y]\times k[y].
\]
Hence the space is disconnected.

\paragraph{What happened mathematically.}
Zero divisors encode reducibility, while idempotents encode a much stronger splitting:
a product of rings.  Geometrically, products of rings correspond to disjoint clopen pieces.
Thus \(V(xy)\) is reducible but connected, whereas two parallel components may be genuinely
disconnected.

\paragraph{Extensions.}
\begin{enumerate}
\item Generalize the result to a decomposition into \(r\) connected components and
orthogonal idempotents \(e_1,\dots,e_r\).
\item Determine all idempotents in \(k^r\).
\item Compare \(k[x,y]/(xy)\) with \(k[x]/(x(x-1))\).
\item Reprove the statement later for \(\operatorname{Spec}A\) without assuming an
algebraically closed base field.
\end{enumerate}

\paragraph{Provenance.}
Adapted from legacy `theory-of-algebraic-geometry-1.tex`, Problem 3.  The scheme-spectrum formulation is deferred; this version is phrased for classical affine algebraic sets.

\fi

